\documentclass[final,5p]{elsarticle}
\usepackage{amsmath,graphicx,url}
\journal{New Astronomy}
\begin{document}
\begin{frontmatter}
\title{A physical ruler for testing printed figure scale}
\author{Layout demonstration}

\begin{abstract}Physical measurements taken from printed figures depend on both the drawing scale and the scale applied when a manuscript is printed. A length ratio can survive a change in page size even when both absolute lengths are wrong. We present a short calibration sheet that places a ten-millimetre segment in one column and compares it with a twenty-millimetre segment in a figure spanning the page. Both figures use the same units and have their own captions. The expected ratio is two at every uniform magnification. The expected absolute lengths are ten and twenty millimetres only when the document is printed at actual size. Measuring the short segment in both positions also allows the reader to distinguish a change affecting one figure from a change applied to the entire sheet. The check requires a ruler and the printed page. It can be repeated after a change of document class, figure width or printer settings.\end{abstract}
\end{frontmatter}
\section{A physical scale for page proofs}
A millimetre ruler makes scaling visible after the paper is printed. Figure~\ref{fig:one} shows a single segment and Fig.~\ref{fig:two} compares two lengths. We follow the SI definition of length \cite{si}.
\begin{equation} L_{\mathrm{mm}}=10 L_{\mathrm{cm}}.\end{equation}
\begin{figure}[ht]\centering\typeout{SINGLE=\the\linewidth}
\rule{10mm}{0.6bp}\caption{One centimetre at print size. The horizontal rule spans ten millimetres.}\label{fig:one}\end{figure}
\begin{figure*}[t]\centering\typeout{DOUBLE=\the\linewidth}
(a) \rule{10mm}{0.6bp}\qquad (b) \rule{20mm}{0.6bp}
\caption{Two centimetres double the reference length. Panels a and b show one and two centimetres on the same physical scale.}\label{fig:two}\end{figure*}
\section{Reading the proof}
Measure the two rules after printing. Supplementary Note~\ref{sec:print} gives the printer setting.
\section{Supplementary Note}\label{sec:print}
Print at actual size. Scaling to fit changes both lengths.
\begin{thebibliography}{1}
\makeatletter
\@ifpackageloaded{natbib}{\bibitem[BIPM(2019)]{si}}{\bibitem{si}}
\makeatother BIPM. The International System of Units (SI), ninth edition, 2019. doi:10.59161/AUEZ1291.
\end{thebibliography}
\end{document}

